% =====================================================================
% 9. Discusión y recomendaciones
% =====================================================================
\section{Discusión y recomendaciones}
\label{sec:discusion}

\subsection{Por qué difieren las fuentes}

Los datos muestran tres causas:

\begin{itemize}
  \item \textbf{Tamaño de las huellas.} Google tiene 1,53 veces más edificios
        pero casi la misma área total, porque divide la misma superficie en
        huellas más pequeñas (82\,m\textsuperscript{2} de media frente a
        126\,m\textsuperscript{2}). Cada modelo recorta de forma distinta los
        techos contiguos. Por eso el \emph{conteo} de edificios no sirve para
        comparar fuentes, y el ranking y la estimación de potencial usan
        \emph{área}.
  \item \textbf{Huecos de detección en el Pacífico.} En el Chocó y la costa
        nariñense Google detecta muy poco o nada. Una causa probable es la
        nubosidad casi permanente de la región, que reduce las imágenes
        útiles para el modelo. Ahí se concentran 18 de los 23 municipios con
        desacuerdo fuerte.
  \item \textbf{Fechas e imágenes distintas.} Microsoft usa imágenes de 2014
        a 2021 y Google una inferencia de 2023; una cabecera que creció en
        esos años tiene más área en la fuente más reciente.
\end{itemize}

La consecuencia práctica es que el total regional es robusto (las fuentes
difieren menos del 0,5\,\% y el umbral lo mueve como mucho un 19\,\%), pero la
cifra de un municipio concreto depende de la fuente, sobre todo en el
Pacífico. Por eso se reporta el área conservadora y, donde las fuentes no
coinciden, conviene reportar un rango y no una sola cifra.

\subsection{Recomendaciones para la UPME}
\label{sec:recomendaciones}

\textbf{Selección de pilotos.}
\begin{enumerate}
  \item Partir del top 15 por área conservadora (tabla~\ref{tab:ranking}):
        concentra el 39\,\% del área de techos de los 170 municipios y ninguno
        tiene las fuentes en desacuerdo.
  \item Dentro del top 15, empezar por los municipios con mayor porcentaje de
        techos, donde el área está concentrada en un casco urbano y no
        dispersa en veredas, lo que facilita la conexión y la operación:
        \textbf{Santa Marta}, \textbf{Santander de Quilichao} y
        \textbf{Apartadó} están en el top 15 y entre los cuatro con mayor
        porcentaje de techos.
  \item Buscar diversidad territorial: el top 15 cubre 11 de las 16
        subregiones, así que un conjunto pequeño de pilotos puede
        representar regiones distintas.
\end{enumerate}

\textbf{Municipios del Pacífico.} No descartarlos por su área baja: en los 23
municipios con desacuerdo fuerte el resultado depende más de la fuente que del
municipio. Se recomienda verificarlos antes de decidir, con imágenes
satelitales recientes o de alta resolución, con una tercera fuente (por
ejemplo GlobalBuildingAtlas, que además trae alturas) o en campo, empezando
por los cinco municipios sin edificios en Google.

\textbf{Refinar la estimación antes de dimensionar.}
\begin{enumerate}
  \item Cruzar el área con la irradiación por municipio del Atlas de
        Radiación Solar~\cite{atlas} en lugar del promedio nacional.
  \item Estimar el área útil por techo (sombras, orientación, inclinación,
        material) con imágenes o modelos de elevación en los municipios
        candidatos; el factor de aprovechamiento de 0,40 es un supuesto.
  \item Cruzar con la demanda eléctrica y la red de cada municipio, para
        dimensionar según el consumo y la capacidad de conexión, y no solo
        según el techo disponible.
\end{enumerate}

\textbf{Continuidad de la plataforma de datos.}
\begin{itemize}
  \item Las fuentes se actualizan: el flujo permite recargar una versión
        nueva de cualquiera de ellas con \campo{make edificios} y recalcular
        todo con \campo{make analisis}, conservando la trazabilidad en
        \campo{fuentes} y \campo{bitacora\_cargas}.
  \item Para extender el análisis a todo el país (unos 16,5 millones de
        edificios solo en Google), el diseño no cambia: la asignación ya
        corre municipio por municipio con índice, y MongoDB permite repartir
        las colecciones de edificios con \emph{sharding} por
        \campo{cod\_mpio}.
\end{itemize}

\subsection{Limitaciones}
\label{sec:limitaciones}

\begin{itemize}
  \item El área es la huella en planta: es un límite superior del área útil
        para paneles.
  \item Las huellas vienen de modelos de aprendizaje automático sobre
        imágenes de fechas distintas; no son un catastro y tienen errores de
        omisión y de partición de techos.
  \item La estimación de potencia usa una irradiación nacional y un factor de
        aprovechamiento fijo.
  \item Un centroide justo sobre un límite puede quedar fuera por la
        diferencia entre aristas esféricas y planas (un edificio de 4,5
        millones en este proyecto).
  \item La validación independiente cubre una muestra de 10 municipios, no
        los 170, y el proyecto no tiene pruebas automáticas; la verificación
        se apoya en controles de conteo en cada etapa.
\end{itemize}
